\documentclass{article}
\usepackage[utf8]{inputenc}
\usepackage{amsmath}
\usepackage{biblatex}
\addbibresource{../bibliography/references_postdoc.bib}

\title{Artificial Childhoods: Critical Learning Windows in RLHF Alignment}
\author{Antigravity Research Lab}
\date{\today}

\begin{document}

\maketitle

\begin{abstract}
Biological brains lose neuroplasticity over time. Do LLMs? This paper investigates the "Critical Window" hypothesis \cite{CriticalPeriodsDL2024}. We find that applying RLHF \textit{after} the crystallization of dense representations (Steps $> 10k$) leads to superficial "Masking" behavior rather than deep alignment.
\end{abstract}

\section{Introduction}
While \textit{Paper II: Causal Editing} demonstrated that we can transiently control model behavior, we encountered a resistance mechanism: the model's "Nature" often asserts itself over our "Nurture" interventions.
This introduces the \textbf{Timing Bottleneck}. We hypothesize that there exists a "Critical Window" during training where the model's fundamental topological structure is malleable. Beyond this point, alignment is merely superficial "masking."

\section{Methodology}
We define \textbf{Alignment Depth (AD)} as the similarity between the base model's probability distribution and the aligned model's distribution in the concept space.
\begin{equation}
    AD = CosineSim(SAE(x_{base}), SAE(x_{aligned}))
\end{equation}

\section{Results}
Our simulations (\texttt{critical\_window\_trainer.py}) reveal:
\begin{itemize}
    \item \textbf{Early:} High Plasticity. The model rewrites its core concepts.
    \item \textbf{Late:} Low Plasticity. The model learns a "Wrapper" or "Persona" around its core (misaligned) concepts.
\end{itemize}

\section{Conclusion}
Safety cannot be an afterthought. There exists a mathematically definable "Childhood" for AI, after which its fundamental worldview is immutable.

\printbibliography

\end{document}
